\usepackage{amsmath, amsthm, amsfonts, amssymb, mathtools, mathrsfs}

\usepackage{iftex}
\usepackage[T1]{fontenc}
% 中文支持: 两个引擎, 两个平台都能直接编译
\ifPDFTeX
  % pdfLaTeX: CJK 宏包 + CJK 环境 (不支持自动中英文间距)
  \usepackage{CJKutf8}
\fi
\ifXeTeX
  % XeLaTeX: xeCJK 提供自动中英文间距
  \usepackage{xeCJK}
  \xeCJKsetup{CJKecglue={\hskip 0.3em plus 0.1em minus 0.1em}}
  % 正文中的 CJK 环境在 xeCJK 下多余, 这里定义为空以免改正文
  \newenvironment{CJK}[2]{}{}
  % 中文字体: 按当前系统已安装的字体自动选择 (Linux / Windows / macOS)
  \IfFontExistsTF{FandolSong-Regular}
    {% TeX Live 自带 Fandol, 各平台均可用
      \setCJKmainfont{FandolSong-Regular}[BoldFont=FandolSong-Bold,ItalicFont=FandolKai-Regular,AutoFakeSlant=0.2]
      \setCJKsansfont{FandolHei-Regular}[BoldFont=FandolHei-Bold]
      \setCJKmonofont{FandolFang-Regular}[AutoFakeSlant=0.2]}
    {%
      \IfFontExistsTF{SimSun}
        {% Windows: 宋体 / 黑体 / 楷体 / 仿宋
          \setCJKmainfont{SimSun}[BoldFont=SimHei,ItalicFont=KaiTi,AutoFakeSlant=0.2]
          \setCJKsansfont{SimHei}[AutoFakeBold=2.5]
          \setCJKmonofont{FangSong}[AutoFakeSlant=0.2]}
        {%
          \IfFontExistsTF{Source Han Serif SC}
            {% Linux: 思源宋体 / 思源黑体
              \setCJKmainfont{Source Han Serif SC}[BoldFont={Source Han Serif SC Bold},ItalicFont={Source Han Serif SC},AutoFakeSlant=0.2]
              \setCJKsansfont{Source Han Sans SC}[BoldFont={Source Han Sans SC Bold},AutoFakeBold=2.5]
              \setCJKmonofont{Source Han Serif SC}[AutoFakeSlant=0.2]}
            {% Linux: Noto CJK / 文鼎
              \IfFontExistsTF{Noto Serif CJK SC}
                {\setCJKmainfont{Noto Serif CJK SC}[AutoFakeBold=2.5,AutoFakeSlant=0.2]
                 \setCJKsansfont{Noto Sans CJK SC}[AutoFakeBold=2.5]
                 \setCJKmonofont{Noto Sans Mono CJK SC}[AutoFakeSlant=0.2]}
                {\setCJKmainfont{AR PL UMing CN}[AutoFakeBold=2.5,AutoFakeSlant=0.2]
                 \setCJKsansfont{AR PL UKai CN}[AutoFakeBold=2.5]
                 \setCJKmonofont{AR PL UKai CN}[AutoFakeSlant=0.2]}}}
    }
\fi
\ifPDFTeX\else\ifXeTeX\else
  \PackageWarningNoLine{MathPreamble}{Only pdfLaTeX and XeLaTeX are supported}
\fi\fi
%\usepackage[utf8]{inputenc}
\usepackage[usenames,dvipsnames,svgnames]{xcolor}
%\usepackage{zhspacing}
\usepackage{tikz}
\usepackage{pstricks}
\usepackage{asymptote}
\usepackage{tikz-dependency}
\usetikzlibrary{calc, datavisualization, cd}
\tikzset{every picture/.append style={remember picture}}
\usepackage{pgfplots}
%\usepackage{yhmath}
\usepackage{fourier}
% zhspacing 只支持 XeTeX, pdfLaTeX 下会报错, 已由上面的 xeCJK 取代
%\usepackage{zhspacing}
\usepackage{tkz-euclide}
%\usetkzobj{all} % for new version of tkz-euclide, this line not needed.
\pgfplotsset{compat=newest}
\usepackage{enumerate}
\usepackage{tcolorbox}
\usepackage[a4paper,hmargin={1cm,1cm},vmargin={2.5cm,2.5cm}]{geometry}
\usepackage{hyperref}
\usepackage{makeidx}
\usepackage{extarrows}
\usepackage{cancel}
\usepackage{microtype}
\usepackage{tikz}
\usepackage{tikz-qtree}
\makeindex

\hypersetup{unicode,colorlinks=true}
%%%%%%%%%%%%%%%%%%%%%%%%%%%%%%%%%%%%%%%%%%%%%%%%%%%%%%%%%%%%%%%%%%%%%%%%%%%%%%%%%%%
%%%%%%%%%%%%%%%%%%%%%%%%%%%%%%%%%%%%%%%%%%%%%%%%%%%%%%%%%%%%%%%%%%%%%%%%%%%%%%%%%%%
%%%%%%%%%%%%%%%%%%%%%%%%%%%%%%%%%%%%%%%%%%%%%%%%%%%%%%%%%%%%%%%%%%%%%%%%%%%%%%%%%%%

% from tcolorbox.pdf
\tcbuselibrary{theorems, 
	breakable, 
	skins, 
	raster, 
	listings, 
	minted,
	xparse, 
	listingsutf8, 
	documentation}
\tcbset{highlight math style={enhanced, colframe=red, colback=red!10!white,boxsep=0pt}}

\newtcbtheorem[number within=section]{question}{问题}%
  {colback=green!5,colframe=green!35!black,fonttitle=\bfseries}{q}
\newtcbtheorem[number within=section]{definition}{定义}%
  {colback=blue!5,colframe=gray!50!black,fonttitle=\bfseries}{def}
\newtcbtheorem[number within=section]{lemma}{引理}%
  {colback=purple!5,colframe=yellow!80!black,fonttitle=\bfseries}{lm}
\newtcbtheorem[number within=section]{theorem}{定理}%
  {colback=purple!5,colframe=red!50!yellow,fonttitle=\bfseries}{th}
\newtcbtheorem[number within=section]{corrolary}{推论}%
  {colback=yellow!5,colframe=yellow!75!black,fonttitle=\bfseries}{co}
\newtcbtheorem[number within=section]{undone}{未知}%
  {colback=red!5,colframe=red!75!black,fonttitle=\bfseries}{un}
\newtcbtheorem[number within=section]{comments}{评论}%
  {colback=gray!5,colframe=gray!75!black,fonttitle=\bfseries}{cm}
\newtcbtheorem[number within=section]{answer}{解答}%
  {colback=gray!5,colframe=gray!75!black,fonttitle=\bfseries}{ans}
% from amsthdoc.pdf
\makeatletter
\renewenvironment{proof}[1][\proofname]{\par
    \pushQED{\qed}%
    \normalfont \topsep6\p@\@plus6\p@\relax
    \trivlist
    \item\relax
    {\itshape
    #1\@addpunct{.}}\hspace\labelsep\ignorespaces
  }{%
    \popQED\endtrivlist\@endpefalse
  }
\makeatother

%%%%%%%%%%%%%%%%%%%%%%%%%%%%%%%%%%%%%%%%%%%%%%%%%%%%%%%%%%%%%%%%%%%%%%%%%%%%%%%%%%%
%%%%%%%%%%%%%%%%%%%%%%%%%%%%%%%%%%%%%%%%%%%%%%%%%%%%%%%%%%%%%%%%%%%%%%%%%%%%%%%%%%%
%%%%%%%%%%%%%%%%%%%%%%%%%%%%%%%%%%%%%%%%%%%%%%%%%%%%%%%%%%%%%%%%%%%%%%%%%%%%%%%%%%%

\def\degree{^\circ}
\def\bq{\begin{question}}
\def\eq{\end{question}}
\def\bqq{\begin{question*}}
\def\eqq{\end{question*}}
\def\bt{\begin{theorem}}
\def\et{\end{theorem}}
\def\bl{\begin{lemma}}
\def\el{\end{lemma}}
\def\bd{\begin{definition}[enhanced jigsaw,breakable]}
\def\ed{\end{definition}}
\def\bc{\begin{corrolary}}
\def\ec{\end{corrolary}}
\def\bm{\begin{comments*}}
\def\em{\end{comments*}}
\def\bu{\begin{undone}}
\def\eu{\end{undone}}
\def\ba{\begin{proof}[解]}
\def\ea{\end{proof}}
\def\bs{\begin{answer*}{}{}}
\def\es{\end{answer*}}
\def\ue{\mathrm{e}}
\def\ud{\,\mathrm{d}}
\def\GF{\mathrm{GF}}
\def\ui{\mathrm{i}}
\def\udiv{\mathrm{div}}
\def\Re{\mathrm{Re}}
\def\Im{\mathrm{Im}}
\def\Res{\mathrm{Res}}
\def\diag{\,\mathrm{diag}\,}
\def\be{\begin{equation}}
\def\ee{\end{equation}}
\def\bee{\begin{equation*}}
\def\eee{\end{equation*}}
\def\sumcyc{\sum\limits_{cyc}}
\def\prodcyc{\prod\limits_{cyc}}
\def\i{\infty}
\def\a{\alpha}
\def\b{\beta}
\def\g{\gamma}
\def\d{\delta}
\def\l{\lambda}
\def\m{\mu}
\def\t{\theta}
\def\kong{\underline{\phantom{asdf}}}

%%%%%%%%%%%%%%%%%%%%%%%%%%%%%%%%%%%%%%%%%%%%%%%%%%%%%%%%%%%%%%%%%%%%%%%%%%%%%%%%%%%
%%%%%%%%%%%%%%%%%%%%%%%%%%%%%%%%%%%%%%%%%%%%%%%%%%%%%%%%%%%%%%%%%%%%%%%%%%%%%%%%%%%
%%%%%%%%%%%%%%%%%%%%%%%%%%%%%%%%%%%%%%%%%%%%%%%%%%%%%%%%%%%%%%%%%%%%%%%%%%%%%%%%%%%

\DeclareMathOperator\arctanh{arctanh}

%%%%%%%%%%%%%%%%%%%%%%%%%%%%%%%%%%%%%%%%%%%%%%%%%%%%%%%%%%%%%%%%%%%%%%%%%%%%%%%%%%%
%%%%%%%%%%%%%%%%%%%%%%%%%%%%%%%%%%%%%%%%%%%%%%%%%%%%%%%%%%%%%%%%%%%%%%%%%%%%%%%%%%%
%%%%%%%%%%%%%%%%%%%%%%%%%%%%%%%%%%%%%%%%%%%%%%%%%%%%%%%%%%%%%%%%%%%%%%%%%%%%%%%%%%%

\providecommand{\abs}[1]{\left\lvert#1\right\rvert}
\providecommand{\norm}[1]{\lVert#1\rVert}
\providecommand{\paren}[1]{\left(#1\right)}
\providecommand{\emph}[1]{\textbf{#1}}

%%%%%%%%%%%%%%%%%%%%%%%%%%%%%%%%%%%%%%%%%%%%%%%%%%%%%%%%%%%%%%%%%%%%%%%%%%%%%%%%%%%
%%%%%%%%%%%%%%%%%%%%%%%%%%%%%%%%%%%%%%%%%%%%%%%%%%%%%%%%%%%%%%%%%%%%%%%%%%%%%%%%%%%
%%%%%%%%%%%%%%%%%%%%%%%%%%%%%%%%%%%%%%%%%%%%%%%%%%%%%%%%%%%%%%%%%%%%%%%%%%%%%%%%%%%

\newcommand{\ZZ}{\mathbb{Z}}
\newcommand{\NN}{\mathbb{N}}
\newcommand{\CC}{\mathbb{C}}
\newcommand{\QQ}{\mathbb{Q}}
\newcommand{\RR}{\mathbb{R}}
\newcommand{\pNN}{\mathbb{N}_{+}}
\newcommand{\notea}[1]{\colorbox{red!20}{#1}}
\newcommand{\noteb}[1]{\tcbhighmath{#1}}